% from ega2-1-fr.tex
\oldpage[II]{104}

\begin{definition}[5.5.2]
\label{II.5.5.2-fr}
On dit qu'un $Y$-préschéma $X$ est \emph{projectif sur $Y$} ou est un
\emph{$Y$-schéma projectif} s'il vérifie les conditions équivalentes a) et b)
de (\hyperref[II.5.5.1-fr]{5.5.1}). On dit qu'un morphisme $f:X\to Y$ est
\emph{projectif} s'il fait de $X$ un $Y$-schéma projectif.
\end{definition}

Il est clair que si $f:X\to Y$ est projectif, il existe un
$\mathcal{O}_X$-Module \emph{très ample relativement à $f$}
(\hyperref[II.4.4.2-fr]{4.4.2}).

\begin{theorem}[5.5.3]
\label{II.5.5.3-fr}
\begin{enumerate}
  \item[\rm{(i)}] Tout morphisme projectif est quasi-projectif et propre.
  \item[\rm{(ii)}] Inversement, soit $Y$ un schéma quasi-compact ou un
  préschéma dont l'espace sous-jacent est noethérien~; alors tout morphisme
  $f:X\to Y$ qui est quasi-projectif et propre est projectif.
\end{enumerate}
\end{theorem}

(i) Il est clair que si $f:X\to Y$ est projectif, il est de type fini et
quasi-projectif (donc en particulier séparé)~; d'autre part, il résulte
aussitôt de (\hyperref[II.5.5.1-fr]{5.5.1, b)}) et de
(\hyperref[II.3.5.3-fr]{3.5.3}) que si $f$ est projectif, il en est de même de
$f\times_Y1_{Y'}:X\times_YY'\to Y'$ pour tout morphisme $Y'\to Y$. Pour
montrer que $f$ est universellement fermé, tout revient donc à prouver qu'un
morphisme projectif $f$ est \emph{fermé}. La question étant locale sur $Y$, on
peut supposer que $Y=\operatorname{Spec}(A)$, donc
(\hyperref[II.5.5.1-fr]{5.5.1}) $X=\operatorname{Proj}(S)$, où $S$ est une
$A$-algèbre graduée engendrée par un nombre fini d'éléments de $S_1$. Pour
tout $y\in Y$, la fibre $f^{-1}(y)$ s'identifie à
$\operatorname{Proj}(S)\times_Y\operatorname{Spec}(k(y))$
(\hyperref[I.3.6.1-fr]{I, 3.6.1}), donc à
$\operatorname{Proj}(S\otimes_Ak(y))$
(\hyperref[II.2.8.10-fr]{2.8.10})~; par suite, $f^{-1}(y)$ est vide si et
seulement si $S\otimes_Ak(y)$ vérifie la condition (TN)
(\hyperref[II.2.7.4-fr]{2.7.4}), autrement dit $S_n\otimes_Ak(y)=0$ pour $n$
assez grand. Mais comme $(S_n)_y$ est un $\mathcal{O}_y$-module de type fini,
la condition précédente signifie que $(S_n)_y=0$ pour $n$ assez grand, en
vertu du lemme de Nakayama. Si $\mathfrak a_n$ est l'annulateur dans $A$ du
$A$-module $S_n$, la condition précédente signifie encore que
$\mathfrak a_n\not\subset\mathfrak j_y$ pour $n$ assez grand
(\hyperref[0.1.7.4-fr]{0, 1.7.4}). Or, comme $S_nS_1=S_{n+1}$ par hypothèse,
on a $\mathfrak a_n\subset\mathfrak a_{n+1}$, et si $\mathfrak a$ est la
réunion des $\mathfrak a_n$, on voit donc que $f(X)=V(\mathfrak a)$, ce qui
prouve que $f(X)$ est fermé dans $Y$. Si maintenant $X'$ est une partie
fermée quelconque de $X$, il existe un sous-préschéma fermé de $X$ ayant $X'$
pour espace sous-jacent (\hyperref[I.5.2.1-fr]{I, 5.2.1}) et il est clair
(\hyperref[II.5.5.1-fr]{5.5.1, a)}) que le morphisme
$X'\to X\xrightarrow{f}Y$ est projectif, donc $f(X')$ est fermé dans $Y$.

(ii) L'hypothèse sur $Y$ et le fait que $f$ est quasi-projectif entraînent
l'existence d'un $\mathcal{O}_Y$-Module quasi-cohérent de type fini
$\mathcal{E}$ et d'une $Y$-immersion $j:X\to\mathbf{P}(\mathcal{E})$
(\hyperref[II.5.3.2-fr]{5.3.2}). Mais comme $f$ est propre, $j$ est fermée par
(\hyperref[II.5.4.4-fr]{5.4.4}), donc $f$ est projectif.

\par\medskip\noindent\textit{Remarques (5.5.4). ---\ }
\label{II.5.5.4-fr}
\begin{enumerate}
  \item[\rm{(i)}] Soit $f:X\to Y$ un morphisme tel que~:
  1\textsuperscript{o} $f$ soit propre~; 2\textsuperscript{o} il existe un
  $\mathcal{O}_X$-Module $\mathcal{L}$ très ample relativement à $f$~;
  3\textsuperscript{o} le $\mathcal{O}_Y$-Module quasi-cohérent
  $\mathcal{E}=f_*(\mathcal{L})$ soit de type fini. Alors $f$ est un morphisme
  projectif~: en effet (\hyperref[II.4.4.4-fr]{4.4.4}), il y a alors une
  $Y$-immersion $r:X\to\mathbf{P}(\mathcal{E})$, et comme $f$ est propre, $r$
  est une immersion fermée (\hyperref[II.5.4.4-fr]{5.4.4}). On verra au
  chapitre III, \S~3, que lorsque $Y$ est localement noethérien, la condition
  3\textsuperscript{o} ci-dessus est une conséquence des deux autres, donc les
  conditions 1\textsuperscript{o} et 2\textsuperscript{o} caractérisent dans
  ce cas les morphismes projectifs et si $Y$ est quasi-compact, on peut
  remplacer la condition 2\textsuperscript{o} par l'hypothèse de l'existence
  d'un $\mathcal{O}_X$-Module ample pour $f$
  (\hyperref[II.4.6.11-fr]{4.6.11}).
  \item[\rm{(ii)}] Soit $Y$ un schéma quasi-compact tel qu'il existe un
  $\mathcal{O}_Y$-Module ample. Pour qu'un $Y$-schéma $X$ soit projectif, il
  faut et il suffit qu'il soit $Y$-isomorphe à un sous-$Y$-schéma fermé d'un
  fibré projectif de la forme $\mathbf{P}_Y^r$. La condition est évidemment
  suffisante.

